% GENERATED FILE. Edit canonical lessons, then run npm run generate:books.
% canonical-source-hash: fnv1a64:b75f65fb04f8e242
% canonical-lessons: RU-C321-korol, RU-C321-koroleva, RU-C321-printsessa, RU-C321-prints, RU-C321-korolevstvo

\chapter{Kings and queens}
\label{ch:ru-kings-and-queens}
\hlchaptermodality{321}

\begin{chapteropening}
\textbf{By the end of this chapter:} \emph{I can talk about a king, a queen, a princess, a prince and a kingdom.}

Complete the last lesson of chapter 321: I can talk about a king, a queen, a princess, a prince and a kingdom.
\end{chapteropening}

\section[\texorpdfstring{\ru{король} (koról)}{король (koról)}]{\ru{король} (koról) — a king}

\label{lesson:RU-C321-korol}



\begin{quote}
Before the new one: say the Russian for a dragon, then the Russian for to turn into.
\end{quote}

\subsection*{You'll want to know: \ru{король}}
\textbf{\ru{король}} (\emph{koról}) — \textquotedblleft{}a king\textquotedblright{}.

\begin{quote}
\ru{Король} \ru{жил} \ru{в} \ru{большом} \ru{замке}.
\end{quote}

\begin{quote}
\emph{Koról zhil v bolshóm zámke.}
\end{quote}

\begin{quote}
The king lived in a big castle.
\end{quote}

Genitive \emph{\ru{короля}}. A Russian ruler was a \emph{\ru{царь}}.

\begin{cousinweb}
\emph{\ru{король}} is inherited through Old East Slavic \emph{korolĭ} from Proto-Slavic \emph{korľĭ} \textquotedblleft{}king\textquotedblright{} (EtymDB), which is usually traced to the name of a great Frankish emperor.
\end{cousinweb}

\begin{grammarlens}[title={What you've built}]
One more word for telling, explaining and describing at B1.
\end{grammarlens}

\subsection*{Your turn}
\begin{itemize}
\raggedright
  \item \emph{Say it:} \emph{koról}
  \item \emph{Say it:} \emph{koról}, once more
  \item \emph{Say it:} a sentence of your own with \emph{koról}
  \item \emph{Recall:} say the Russian for a dragon, then the Russian for to turn into, then say \emph{koról} again
\end{itemize}

\begin{tcolorbox}[breakable,colback=teal!4,colframe=teal!35!black,title={Before you move on}]
What does \ru{король} mean? (\textquotedblleft{}A king\textquotedblright{}.) Say it once more.
\end{tcolorbox}
\section[\texorpdfstring{\ru{королева} (koroléva)}{королева (koroléva)}]{\ru{королева} (koroléva) — a queen}

\label{lesson:RU-C321-koroleva}



\begin{quote}
Before the new one: say the Russian for to turn into, then the Russian for a king.
\end{quote}

\subsection*{You'll want to know: \ru{королева}}
\textbf{\ru{королева}} (\emph{koroléva}) — \textquotedblleft{}a queen\textquotedblright{}.

\begin{quote}
\ru{Королева} \ru{улыбалась} \ru{и} \ru{махала} \ru{рукой}.
\end{quote}

\begin{quote}
\emph{Koroléva ulybálas i makhála rukóy.}
\end{quote}

\begin{quote}
The queen smiled and waved.
\end{quote}

Feminine.

\begin{cousinweb}
\emph{\ru{королева}} is built on \emph{\ru{король}} \textquotedblleft{}king\textquotedblright{}, inherited from Proto-Slavic \emph{korľĭ} (EtymDB), with the feminine ending \emph{-\ru{ева}}.
\end{cousinweb}

\begin{grammarlens}[title={What you've built}]
One more word for telling, explaining and describing at B1.
\end{grammarlens}

\subsection*{Your turn}
\begin{itemize}
\raggedright
  \item \emph{Say it:} \emph{koroléva}
  \item \emph{Say it:} \emph{koroléva}, once more
  \item \emph{Say it:} a sentence of your own with \emph{koroléva}
  \item \emph{Recall:} say the Russian for to turn into, then the Russian for a king, then say \emph{koroléva} again
\end{itemize}

\begin{tcolorbox}[breakable,colback=teal!4,colframe=teal!35!black,title={Before you move on}]
What does \ru{королева} mean? (\textquotedblleft{}A queen\textquotedblright{}.) Say it once more.
\end{tcolorbox}
\section[\texorpdfstring{\ru{принцесса} (printséssa)}{принцесса (printséssa)}]{\ru{принцесса} (printséssa) — a princess}

\label{lesson:RU-C321-printsessa}



\begin{quote}
Before the new one: say the Russian for a king, then the Russian for a queen.
\end{quote}

\subsection*{You'll want to know: \ru{принцесса}}
\textbf{\ru{принцесса}} (\emph{printséssa}) — \textquotedblleft{}a princess\textquotedblright{}.

\begin{quote}
\ru{Принцесса} \ru{не} \ru{хотела} \ru{выходить} \ru{замуж}.
\end{quote}

\begin{quote}
\emph{Printséssa ne khotéla vykhodít zámuzh.}
\end{quote}

\begin{quote}
The princess did not want to marry.
\end{quote}

Feminine.

\begin{cousinweb}
\emph{\ru{принцесса}} was borrowed from French \emph{princesse} (EtymDB), built on \emph{prince}, from Latin \emph{prīnceps} \textquotedblleft{}first, chief\textquotedblright{}.
\end{cousinweb}

\begin{grammarlens}[title={What you've built}]
One more word for telling, explaining and describing at B1.
\end{grammarlens}

\subsection*{Your turn}
\begin{itemize}
\raggedright
  \item \emph{Say it:} \emph{printséssa}
  \item \emph{Say it:} \emph{printséssa}, once more
  \item \emph{Say it:} a sentence of your own with \emph{printséssa}
  \item \emph{Recall:} say the Russian for a king, then the Russian for a queen, then say \emph{printséssa} again
\end{itemize}

\begin{tcolorbox}[breakable,colback=teal!4,colframe=teal!35!black,title={Before you move on}]
What does \ru{принцесса} mean? (\textquotedblleft{}A princess\textquotedblright{}.) Say it once more.
\end{tcolorbox}
\section[\texorpdfstring{\ru{принц} (prints)}{принц (prints)}]{\ru{принц} (prints) — a prince}

\label{lesson:RU-C321-prints}



\begin{quote}
Before the new one: say the Russian for a queen, then the Russian for a princess.
\end{quote}

\subsection*{You'll want to know: \ru{принц}}
\textbf{\ru{принц}} (\emph{prints}) — \textquotedblleft{}a prince\textquotedblright{}.

\begin{quote}
\ru{Младший} \ru{принц} \ru{любил} \ru{охоту}.
\end{quote}

\begin{quote}
\emph{Mládshiy prints lyubíl okhótu.}
\end{quote}

\begin{quote}
The younger prince loved hunting.
\end{quote}

A king's son.

\begin{cousinweb}
\emph{\ru{принц}} was borrowed from German \emph{Prinz} (EtymDB), which comes from French \emph{prince}, from Latin \emph{prīnceps} \textquotedblleft{}first, chief\textquotedblright{}, built on \emph{prīmus} \textquotedblleft{}first\textquotedblright{} and \emph{capere} \textquotedblleft{}to take\textquotedblright{}.
\end{cousinweb}

\begin{grammarlens}[title={What you've built}]
One more word for telling, explaining and describing at B1.
\end{grammarlens}

\subsection*{Your turn}
\begin{itemize}
\raggedright
  \item \emph{Say it:} \emph{prints}
  \item \emph{Say it:} \emph{prints}, once more
  \item \emph{Say it:} a sentence of your own with \emph{prints}
  \item \emph{Recall:} say the Russian for a queen, then the Russian for a princess, then say \emph{prints} again
\end{itemize}

\begin{tcolorbox}[breakable,colback=teal!4,colframe=teal!35!black,title={Before you move on}]
What does \ru{принц} mean? (\textquotedblleft{}A prince\textquotedblright{}.) Say it once more.
\end{tcolorbox}
\section[\texorpdfstring{\ru{королевство} (korolévstvo)}{королевство (korolévstvo)}]{\ru{королевство} (korolévstvo) — a kingdom}

\label{lesson:RU-C321-korolevstvo}



\begin{quote}
Before the new one: say the Russian for a princess, then the Russian for a prince.
\end{quote}

\subsection*{You'll want to know: \ru{королевство}}
\textbf{\ru{королевство}} (\emph{korolévstvo}) — \textquotedblleft{}a kingdom\textquotedblright{}.

\begin{quote}
\ru{В} \ru{маленьком} \ru{королевстве} \ru{все} \ru{знали} \ru{друг} \ru{друга}.
\end{quote}

\begin{quote}
\emph{V málenkom korolévstve vse ználi drug drúga.}
\end{quote}

\begin{quote}
In the little kingdom everyone knew each other.
\end{quote}

Neuter.

\begin{cousinweb}
\emph{\ru{королевство}} is built on \emph{\ru{король}} \textquotedblleft{}king\textquotedblright{}, inherited from Proto-Slavic \emph{korľĭ} (EtymDB), with the ending \emph{-\ru{ство}}.
\end{cousinweb}

\begin{grammarlens}[title={What you've built}]
One more word for telling, explaining and describing at B1.
\end{grammarlens}

\subsection*{Your turn}
\begin{itemize}
\raggedright
  \item \emph{Say it:} \emph{korolévstvo}
  \item \emph{Say it:} \emph{korolévstvo}, once more
  \item \emph{Say it:} a sentence of your own with \emph{korolévstvo}
  \item \emph{Recall:} say the Russian for a princess, then the Russian for a prince, then say \emph{korolévstvo} again
\end{itemize}

\begin{tcolorbox}[breakable,colback=teal!4,colframe=teal!35!black,title={Before you move on}]
What does \ru{королевство} mean? (\textquotedblleft{}A kingdom\textquotedblright{}.) Say it once more.
\end{tcolorbox}
